\paragraph{Native representation and boundary state}

We represent motion directly in the robot's configuration space, including full pelvis orientation. A fixed yaw anchor $R_A$ is removed once for a motion sequence, and every frame stores absolute orientation in that anchor:
\begin{align}
 R_t^A&=R_A^\top R_t^{\rm world},\\
 x_t&=[\Delta p_{xy,t}^{\rm local},\;h_t,\;\rho(R_t^A),\;\theta_t].
\label{eq:motion}
\end{align}
Here $\Delta p_{xy,t}^{\rm local}$ denotes planar displacement, $h_t$ is pelvis height, $\rho$ is the six-dimensional rotation representation~\cite{zhou2019rotation}, and $\theta_t$ contains 29 joint angles, giving 38 motion dimensions. This retains roll and pitch without multiplying a sequence of predicted rotation increments. Planar displacements are integrated using the appropriate preceding heading. The anchor remains fixed over the generated sequence; it is not reset at every commit.

The corresponding 71-dimensional boundary state is
\begin{equation}
 b_t=[h_t,v^{\rm local}_{xy,t},v_{z,t},\rho(R_t^A),\omega_t,
 \theta_t,\dot\theta_t].
\label{eq:state}
\end{equation}
The state contains pelvis height, local planar and vertical velocities, orientation, angular velocity $\omega_t$, joint angles, and joint velocities. Conditioning on these quantities makes the starting pose and motion explicit at each segment boundary. They are reconstructed from the committed reference.


\paragraph{Streaming configuration}
The motion rate is 30\,Hz, with two frames per latent step. We use 64 history
steps (128 frames, approximately 4.27\,s), an eight-step plan (16 frames,
approximately 0.53\,s), and a four-step commit (eight frames, approximately
0.27\,s). The codec latent has 16 dimensions and its codebook has 512 entries.
The music condition uses 14 native forecast states at 25\,Hz, spanning
approximately 0.56\,s. Each state retains its physical timestamp.

\paragraph{Reconstruction gradients and auxiliary objectives}
Let $\mathrm{sg}$ denote stop-gradient. For encoder output $z$, embedding $e$, and $r=z-\mathrm{sg}(e)$, we form full and structure-only reconstructions, together with an auxiliary residual-corrupted reconstruction:
\begin{align}
 x_f&=D_\psi(\mathrm{sg}(e)+r,b),\\
 x_s&=D_\psi(z+\mathrm{sg}(e-z),b),\\
 x_c&=D_\psi(\mathrm{sg}(e)+\mathcal C(r),b).
\end{align}
The full branch is supervised on $\Phi_f$; the other two are supervised on $\Phi_s$. The corruption $\mathcal C$ either zeros detail or adds small channel-scaled noise, encouraging the structural reconstruction to persist under changes in the residual. The combined objective is
\begin{equation}
\begin{split}
\mathcal L_{\rm codec}={}&\mathcal L_f+\mathcal L_s
 +0.25\mathcal L_{\rm robust}+\mathcal L_{\rm VQ}\\
 &+\mathcal L_{\rm boundary}.
\end{split}
\label{eq:codec}
\end{equation}
The final term supervises boundary-transition derivatives and pelvis orientation at the beginning of the reconstructed segment. Appendix~\ref{app:training} provides the training details.


\paragraph{Conditioning and sampling}
For forecast state $u_j$ with lead $\ell_j$ relative to the motion decision,
we form $c_j=W\,\mathrm{LN}(u_j)+\phi(\ell_j)$, where $W$ is a learned
projection, $\mathrm{LN}$ denotes layer normalization, and $\phi$ represents
physical time. Timestamp-based masking excludes expired forecast states from
their aligned motion positions. A null condition bypasses music modulation.
The structural planner uses a classification objective and the detail model
a denoising objective. Residual generation uses ten DDIM
steps~\cite{song2021ddim}.

